\begin{textsample}{NEG Dim 5 – human – Score -10.00 – es/es\_ef\_anos\_iniciais\_matematica\_3\_p9.txt}  \label{ex:f5_neg_007}
Os \textbf{verbos} presentes nas \textbf{habilidades} da \textbf{Base} \textbf{Nacional} Comum \textbf{Curricular} e no Currículo do Espírito Santo \textbf{explicitam} os processos \textbf{cognitivos} que se \textbf{espera} sejam desenvolvidos pelos \textbf{estudantes} no processo de ensino e aprendizagem ( KRATHWOHL e ANDERSON, 2001 ).
Dessa forma, a progressão das aprendizagens, que se \textbf{explicita} na comparação das \textbf{habilidades} em cada \textbf{ano}, ou de um \textbf{ano} para o outro, pode estar \textbf{relacionada} aos :

% matched lemmas: ano, base, cognitivo, curricular, esperar, estudante, explicitar, habilidade, nacional, relacionar, verbo
\end{textsample}
